% ============================================================
% technical_specification/requirements_table.tex — functional requirements
% table (SKELETON; filled in by the student).
%
% The source of requirement definitions (\req) for the traceability
% panel. Included from technical_specification.tex
% (\input{requirements_table}) and test_program_and_methodology.tex
% (\input{../technical_specification/requirements_table}).
%
% Each requirement is declared with \req{ID}{Text} in the
% "Requirement" cell: only the text is printed, while the
% "Requirements traceability" panel links the definition to
% \reqref{ID} references in the thesis/TPM/manuals. The ID in the
% first cell and in \req must match; keep the single ТЗ-Ф-* prefix —
% the sample \reqref references in other documents use it.
% ============================================================
\begin{longtable}{|l|>{\raggedright\arraybackslash}p{0.28\textwidth}|>{\raggedright\arraybackslash}p{0.28\textwidth}|>{\centering\arraybackslash}p{0.22\textwidth}|}
\caption{Functional requirements of the system} \label{tab:requirements} \\
\hline
\textbf{No.} & \textbf{User story} & \textbf{Requirement} & \textbf{Scope} \\
\hline
\endfirsthead
\multicolumn{4}{l}{Table \thetable{} --- continued} \\
\hline
\textbf{No.} & \textbf{User story} & \textbf{Requirement} & \textbf{Scope} \\
\hline
\endhead
\hline
\endfoot
\hline
\endlastfoot
ТЗ-Ф-01 & As a user, I want \ldots, so that \ldots & \req{ТЗ-Ф-01}{The system shall …} & \ldots \\
\hline
ТЗ-Ф-02 & As a user, I want \ldots, so that \ldots & \req{ТЗ-Ф-02}{The system shall …} & \ldots \\
\hline
\end{longtable}
